\section*{A4. Short-answer key}
Award up to 5 points per answer. Respect the three-sentence limit: grade the argument, not literary style. Do not require students to use behavioural terminology not introduced in the course.
\subsection*{1. Paying to resist temptation}
\emph{Example answer (three sentences).} With fixed preferences and costless control over choices, removing an option for a fee cannot improve the best attainable outcome. Paying can nevertheless make sense if the app reduces attention costs or commits a future self whose choices differ from today's plans. The explanation must specify which feature is missing from the simple static model rather than treating every restriction as beneficial.

Marking: 2 points for the fixed-preference/costless-choice benchmark, 2 for a coherent attention or commitment explanation and 1 for explicitly identifying the changed assumption. A mere claim that ``everyone has different tastes'' is insufficient. If the explanation values the app's effect on concentration directly, it must make clear that the product then includes a different environment, not merely a smaller set of otherwise identical options.
\subsection*{2. Too cheap to be desirable}
\emph{Example answer (three sentences).} The shopper may value exclusivity or publicly recognised status, so a widely available bag need not be the same fully specified consumption alternative. Price may also provide information about quality. The statement alone does not establish incomplete or intransitive preferences; a meaningful comparison must hold relevant characteristics, information and social context fixed.

Marking: 2 points for a coherent characteristic or information explanation, 2 for stating what must be held fixed and 1 for avoiding an unsupported inference of irrationality. An unexplained label such as ``Veblen good'' is not a complete answer. Both price and availability change in the question, so do not infer a pure own-price effect from it.
